% O2. Where beta <= 1/3 enters.
% Part 1: Section 2.3, the sentence that ends at sections/02_model.tex:123 and its footnote; the footnote text stands in
% place of sections/02_model.tex:123-125.
% Part 2: Appendix A.2, in place of appendix/A2_beliefs_trading.tex:232-236, after the proof of Lemma A.3 (and after
% Lemmas A.4 and A.5 once they are placed), followed by the paragraph on the exclusion bound at beta > 1/3.
% source: A2:57-91 (Lemma A.1 and its proof); A3:51, A3:122; A5:23, A5:97-98, A5:358, A5:376-382; A6:103-104,
% A6:306-308, A6:402-404; A8:110, A8:121, A8:162; referee report W3, "Shared Candidate 1", section 1 (the extension
% of the research results to beta < 1/2); scan/reduction_A3_A5.md, part 2 (the reduction at 0 < beta < 1/2);
% PLAN.md Part 5, O2.

% Part 1 (main text).
\noindent\ldots\ sells, abstains, or buys one unit with probabilities \(\beta\), \(1-2\beta\), and \(\beta\), where
\(0<\beta\le1/3\).\footnote{Only the posterior bound at the zero order uses \(\beta\le1/3\). Through it the bound
enters the exclusion bound of Section~\ref{sec:inference} and the result on preparation cost laws
(Proposition~\ref{prop:atomless}), and the research parameter region of Section~\ref{sec:research} and the
identification of the informed share in Section~\ref{sec:empirical} range over the slots with \(\beta\le1/3\).
Every other result holds for every \(\beta<1/2\) (Appendix~\ref{app:beliefs}).}

% Part 2 (Appendix A.2).
\paragraph{Where \(\beta\le1/3\) enters.}
The bound \(\beta\le1/3\) is the inequality \(b_0\ge b\), and the paper uses it only through the posterior bound at
the zero order, Lemma~\ref{lem:beliefs}(c). That bound enters Corollary~\ref{cor:exclusion}, the zero order in
Lemma~\ref{lem:a5-bounds}(a), and Step~1 of the proof of Proposition~\ref{prop:atomless}, at a price that pools the
zero order. Proposition~\ref{prop:researchregion} and parts (a) and (c) of Proposition~\ref{prop:informedshare}
describe sets of slots, and they range over the slots with \(\beta\le1/3\). Every other result holds for every
\(\beta<1/2\), for three reasons. The bounds \eqref{eq:a2-bounds}, and every argument that forces orders under every
candidate, use the posteriors of the nonzero orders only, which Lemma~\ref{lem:beliefs}(b) bounds for every
\(\beta<1/2\). Every equilibrium tuple computes its posteriors from its stated orders by Bayes' rule, as
Lemma~\ref{lem:beliefs}(d) does for full orders. Lemmas~\ref{lem:prices}, \ref{lem:trading}, \ref{lem:reduction}
and~\ref{lem:twonumbers} and the paragraph on constructed tuples use Lemma~\ref{lem:beliefs}(a) and the conditions
they list.

\paragraph{The exclusion bound at \(\beta>1/3\).}
The exclusion bound needs \(\beta\le1/3\). Let \(1/3<\beta<1/2\), fix \(\lambda,\rho\in(0,1)\) and an auction with
\(\DT>0\) and \(g_H>g_L\), and let class \(H\) abstain and class \(L\) sell. On the orders \((-1,0,1)\),
\(f_H=(b,\,b_0+\lambda,\,b)\) and \(f_L=(b+\lambda,\,b_0,\,b)\), so the posteriors are \(\mu(-1)=m(\pi)\),
\(\mu(1)=\pi\), and \(\mu(0)>M(\pi)\), because \(b_0<b\) gives \(1+\lambda/b_0>1+u\). Take a prior close enough to
one that \(1-\pi<m(\pi)\), a high cost with \(M(\pi)<\tau<\mu(0)\), and a trading cost with
\(\rho\DT(1-\pi)<k<\rho\DT\,m(\pi)\). Entry is \(\rho\) after a sale and after a purchase and one after no order. It
is nondecreasing in the posterior, so the paragraph on constructed tuples gives three distinct prices and the
beliefs at them, and the high-cost challenger prepares exactly at the price of no order, with no tie. Against these
prices the purchase earns \(\rho\DT(1-\pi)-k<0\), the sale earns \(\rho\DT\,m(\pi)-k>0\), and each wrong-signed order
earns at most \(-k\), so abstention is the unique best order of class \(H\) and the sale that of class \(L\). With
the prices of the constructed tuple and truthful bids, this is an equilibrium under all four tie rules in which
\(B(M(\pi))<c_H\) and a high-cost challenger prepares after no order.
